\documentclass{article}
\usepackage[T1]{fontenc}
\usepackage{lmodern}
\usepackage{graphicx}
\usepackage{amsmath,amssymb}
\usepackage[a4paper,margin=2cm]{geometry}
\usepackage[absolute,overlay]{textpos}
\setlength{\TPHorizModule}{1mm}
\setlength{\TPVertModule}{1mm}
\usepackage[indonesian]{babel}
\usepackage{hyperref}
\setlength{\emergencystretch}{2em}
% unit: Halaman 17

\begin{document}

% === Halaman 17 (llm) ===

\begin{itemize}
    \item Jumlah entri (baris) di dalam setiap buku kode untuk adalah $2^n$.
    \item Namun, semakin besar ukuran blok, semakin besar pula ukuran buku kodenya.
    \item Misalkan jika blok berukuran 64 bit, maka buku kode terdiri dari $2^{64}$ buah kode (\textit{entry}), yang berarti terlalu besar untuk disimpan. Lagipula, setiap kunci mempunyai buku kode yang berbeda.
\end{itemize}

\end{document}
